% Part: lambda-calculus
% Chapter: syntax
% Section: de-bruijn

\documentclass[../../../include/open-logic-section]{subfiles}

\begin{document}

\olfileid{lam}{syn}{deb}

\olsection{德布鲁因指标}

α-等价非常自然，因为α-等价的项“含义相同”。事实上，可以为λ-项提供一种语法，使得能够相互α-转换的项在这种语法下不被区分。最著名的方法是用变元的\emph{德布鲁因指标}取代变元名。

当我们写下 $\lambd[x][M]$ 时，便明确指出了 $x$ 是该函数的参数，从而可以在 $M$ 中用 $x$ 来指称该参数。然而在德布鲁因指标中，参数没有名称，在函数体内对它们的指称由一个数字表示，该数字代表参数与其引用处之间相隔的抽象层数。例如，考虑 $\lambd[x][\lambd[y][y x]]$ 的例子：外层抽象绑定的变元是 $x$；内层抽象绑定的变元是 $y$；子项 $y x$ 处于内层抽象的辖域内：$y$ 与其抽象算子 $\lambd[y]$ 之间没有相隔的抽象，而 $x$ 与其抽象算子 $\lambd[x]$ 之间相隔一层抽象。因此，我们用 $0\, 1$ 表示 $y x$，用 $\lambd[][\lambd[][01]]$ 表示整个项。

\begin{defn}
  德布鲁因项归纳定义如下：
  \begin{enumerate}
  \item $n$，其中 $n$ 为任意自然数。
  \item $PQ$，其中 $P$、$Q$ 均为德布鲁因项。
  \item $\lambd[][N]$，其中 $N$ 为德布鲁因项。
  \end{enumerate}
\end{defn}

从普通λ-项到德布鲁因指标项的形式化转换如下：

\begin{defn}
  \begin{align*}
    F_\Gamma(x) &= \Gamma(x) \\
    F_\Gamma(PQ) &= F_\Gamma(P)F_\Gamma(Q) \\
    F_\Gamma(\lambd[x][N]) &= \lambd[][F_{x,\Gamma}(N)]
  \end{align*}
  其中 $\Gamma$ 是一个从 $0$ 开始编号的变元列表，$\Gamma(x)$ 表示变元 $x$ 在 $\Gamma$ 中的位置。
  例如，若 $\Gamma$ 为 $x,y,z$，则 $\Gamma(x)$ 为 $0$，$\Gamma(z)$ 为 $2$。

  $x,\Gamma$ 表示将 $x$ 添加到 $\Gamma$ 首部得到的列表；例如，沿用上例的 $\Gamma$，$w,\Gamma$ 即为 $w,x,y,z$。
\end{defn}

从德布鲁因项恢复标准λ-项的方法如下：

\begin{defn}
  \begin{align*}
    G_\Gamma(n) &= \Gamma[n] \\
    G_\Gamma(PQ) &= G_\Gamma(P) G_\Gamma(Q) \\
    G_\Gamma(\lambd[][N]) &= \lambd[x][G_{x,\Gamma}(N)]
  \end{align*}
  其中 $\Gamma$ 同样是一个从 $0$ 开始编号的变元列表，$\Gamma[n]$ 表示位于位置~$n$ 的变元。例如，
  若 $\Gamma$ 为 $x,y,z$，则 $\Gamma[1]$ 为 $y$。

  最后一个等式中的变元 $x$ 可选为任意不在 $\Gamma$ 中的变元。
\end{defn}

以下给出若干结论，不加证明：

\begin{prop}
  若 $M \aconvone M'$，且 $\Gamma$ 为任意包含 $\FV{M}$ 的列表，则
  $F_\Gamma(M) \eqs F_\Gamma(M')$。
\end{prop}

\end{document}